\documentclass[11pt]{article}
\usepackage[margin=1in]{geometry}
\usepackage{times}
\usepackage{booktabs}
\usepackage{graphicx}
\usepackage{hyperref}
\usepackage{enumitem}
\usepackage{amsmath}

\begin{document}

\section{Introduction}
DNA methylation patterns change systematically with age, making them powerful biomarkers for biological age prediction. However, methylation data poses challenges: high dimensionality (thousands of CpG sites), complex inter-site correlations, and continuous values in [0,1]. Traditional approaches use linear models (e.g., ElasticNet), but recent work shows that Transformer-based models can capture non-linear dependencies between CpG sites.

We base our work on MethylGPT~[1], a foundation model that treats CpG sites as tokens and applies self-attention to learn methylation representations. We propose adapting the BMFM-RNA architecture~[2], originally designed for single-cell RNA-seq, to methylation data. Our hypothesis is that BMFM-RNA's multi-field tokenization---which separately encodes site identities and continuous values---may better preserve fine-grained methylation information compared to MethylGPT's approach.

\begin{itemize}[leftmargin=1.5em]
\item \textbf{Base paper:} \href{https://www.biorxiv.org/content/10.1101/2024.10.30.621013v2}{\textbf{MethylGPT (bioRxiv 2024)}}
\item \textbf{Code:} \href{https://github.com/netanelazran11/BiomedSciAI_ACL_Project}{\textbf{GitHub Repository}}
\end{itemize}

\section{Data Description}
\begin{table}[h]
\centering
\caption{Dataset summary and preprocessing.}
\label{tab:data}
\begin{tabular}{ll}
\toprule
Item & Description \\
\midrule
Pretraining sources & EWAS Data Hub; ClockBase \\
Collected profiles & 226,555 \\
After QC + deduplication & 154,063 \\
Platforms & Illumina 27K / 450K / EPIC \\
CpG harmonization & 49,156 sites ($\geq$5 EWAS traits and/or $\geq$95\% presence) \\
Value range & $\beta$-values in [0,1] \\
Fine-tuning samples & 11,453 (ages 0--100, multi-tissue) \\
\bottomrule
\end{tabular}
\end{table}

\section{Baseline Architecture (MethylGPT)}

\begin{table}[h]
\centering
\caption{MethylGPT architecture (pretraining vs fine-tuning).}
\label{tab:arch_compare}
\begin{tabular}{lll}
\toprule
Component & Pretraining & Fine-tuning \\
\midrule
Base model & TransformerModel (scGPT) & Same (frozen then unfrozen) \\
Transformer & 6 layers, 4 heads, $d=64$ & Same (loaded from pretrained) \\
Task head & ExprDecoder (MLM) & ResNet1D age head  \\
Output & Methylation reconstruction & Single age prediction \\
\bottomrule
\end{tabular}
\end{table}
\subsection{Pretraining}

We pretrain using MLM on 154,063 methylation profiles. The encoder achieves Test PCC = 0.997, indicating accurate reconstruction of masked $\beta$-values.

\section{Loss Function}
Pretraining uses masked MSE to reconstruct methylation values with a 30\% mask ratio. Fine-tuning uses MSE on normalized chronological age.

\section{Evaluation}
We use a fixed 48/12/40 train/val/test split with model selection on validation MAE. Generalization is assessed using 5 random seeds (40--44).

\textbf{MethylGPT fine-tuning:} For age prediction, MethylGPT replaces the MLM head with a ResNet1D regression head using 1$\times$1 convolutions over the token embeddings, followed by adaptive pooling and a linear layer to produce a single age prediction.

\textbf{Metrics:} We report MAE (interpretable in years), MedAE (robust to outliers), and $R^2$ (proportion of variance explained). These are standard metrics for age regression that allow comparison with prior work~[3, 4].

\section{Baseline Results (MethylGPT)}
\begin{table}[h]
\centering
\caption{MethylGPT baseline results (5 seeds).}
\label{tab:results}
\begin{tabular}{lcc}
\toprule
Metric & Mean $\pm$ Std & Best \\
\midrule
Validation MAE (years) & 5.12 $\pm$ 0.89 & 4.13 \\
Test MAE (years) & 4.95 & -- \\
Test MedAE (years) & 3.95 $\pm$ 0.20 & -- \\
Test $R^2$ & 0.911 & -- \\
\bottomrule
\end{tabular}
\end{table}

\section{Proposed Changes: BMFM-RNA Architecture}

We adapt BMFM-RNA~[2], a BERT-based foundation model for RNA-seq ($\sim$110M parameters), to methylation data.

\subsection{Why BMFM-RNA for Methylation?}

\textbf{Multi-field tokenization.} Unlike MethylGPT which combines site and value information implicitly, BMFM-RNA uses explicit dual-field representation:
\begin{equation}
\mathbf{h}_i = \mathbf{e}_{\text{cpg}}(s_i) + \mathbf{e}_{\beta}(\beta_i)
\end{equation}
where $\mathbf{e}_{\text{cpg}}$ captures site-specific information (genomic context) and $\mathbf{e}_{\beta}$ encodes the methylation level via an MLP. This separation allows the model to learn which sites are informative independently of their current value.

\textbf{Continuous value encoding.} The $\beta$-value encoder projects continuous values into embedding space without discretization, preserving fine-grained information.

\textbf{Larger capacity.} BMFM-RNA uses 512-dimensional embeddings and 8 attention heads (vs. MethylGPT's 64-dim, 4 heads), potentially capturing more complex CpG interactions.

\subsection{Architecture Comparison}

\begin{table}[h]
\centering
\caption{Architecture comparison.}
\label{tab:scbert_config}
\begin{tabular}{lll}
\toprule
Parameter & BMFM-RNA (ours) & MethylGPT \\
\midrule
Layers & 6 & 6 \\
Attention heads & 8 & 4 \\
Hidden size $d$ & 512 & 64 \\
Parameters & $\sim$23M & $\sim$2M \\
\bottomrule
\end{tabular}
\end{table}

\subsection{Pretraining on Methylation Data}

We pretrain the adapted BMFM-RNA encoder using MLM on 154,063 methylation profiles for 250 epochs. The model learns to reconstruct masked $\beta$-values from context, achieving Test PCC = 0.997 and Test MAE = 0.0195 on the [0,1] scale. This high reconstruction accuracy indicates the encoder captures meaningful inter-CpG dependencies that transfer to downstream tasks.

\subsection{Fine-Tuning Pipeline}
\begin{itemize}[leftmargin=1.5em]
\item \textbf{Pooling:} Mean pooling over tokens (MLM doesn't train [CLS])
\item \textbf{Age head:} MLP (512 $\to$ 256 $\to$ 128 $\to$ 1) with Dropout(0.2)
\item \textbf{Freeze strategy:} Encoder frozen epochs 0--4, unfrozen epoch 5+
\item \textbf{Normalization:} $z$-score on ages for stable gradients
\end{itemize}

\section{Experimental Setup}

\begin{table}[h]
\centering
\caption{Training hyperparameters.}
\label{tab:training_config}
\begin{tabular}{ll}
\toprule
Hyperparameter & Value \\
\midrule
Optimizer & AdamW (lr=$5 \times 10^{-4}$, wd=0.01) \\
LR schedule & Cosine decay with warmup \\
Batch size & 32 (effective 64) \\
Early stopping & Patience = 60 on val/MAE \\
\bottomrule
\end{tabular}
\end{table}

\section{Results}

\begin{table}[h]
\centering
\caption{Comparison on test set.}
\label{tab:all_comparison}
\begin{tabular}{lcc}
\toprule
Model & Test MAE (years) & Test $R^2$ \\
\midrule
Mean prediction & 22.82 & 0.00 \\
MethylGPT & 4.95 & 0.911 \\
\textbf{BMFM-RNA (ours)} & \textbf{4.85} & \textbf{0.923} \\
\bottomrule
\end{tabular}
\end{table}

BMFM-RNA achieves 2\% improvement in MAE (4.85 vs 4.95 years) and 1.2\% improvement in $R^2$ (0.923 vs 0.911). Both models reduce error by $\sim$78\% compared to mean prediction.

\section{Discussion and Conclusions}

\textbf{Performance.} BMFM-RNA's multi-field tokenization and larger capacity yield modest but consistent improvements over MethylGPT. The explicit separation of site identity and methylation value appears beneficial.

\textbf{Overfitting.} Train MAE = 2.21 years vs test MAE = 4.85 years indicates some overfitting. Dropout(0.2) helped reduce this gap from earlier experiments.

\textbf{Capacity vs. gains.} Despite 10$\times$ more parameters (23M vs 2M), improvement is modest (0.10 years), suggesting methylation age prediction may be approaching a performance ceiling with current approaches.

\textbf{Limitations.} Results are from a single seed for BMFM-RNA; multi-seed experiments would strengthen conclusions.

\section*{References}
\begin{enumerate}[leftmargin=1.5em]
\item Ying, K., et al., ``MethylGPT: a foundation model for the DNA methylome,'' \emph{bioRxiv}, 2024.
\item Dandala, B., et al., ``BMFM-RNA: An Open Framework for Building Transcriptomic Foundation Models,'' \emph{arXiv}, 2025.
\item Horvath, S., ``DNA methylation age of human tissues and cell types,'' \emph{Genome Biology}, 14(10):R115, 2013.
\item Hannum, G., et al., ``Genome-wide methylation profiles reveal quantitative views of human aging rates,'' \emph{Molecular Cell}, 49(2):359--367, 2013.
\item de Lima Camillo, L.P., et al., ``AltumAge: A pan-tissue DNA methylation epigenetic clock,'' \emph{npj Aging}, 8(1):1--15, 2022.
\end{enumerate}

\end{document}
